\subsection{Construction of the port graphs}\label{app:refresh-graph}

The proof of Proposition~\ref{lem:refresh-graph} starts on a finite
Grassmannian. The quadrilateral inequality gives an upper bound on the
mean squared distance between complementary subspaces in terms of the
mean for adjacent ones. Equality would realize the excluded cube, so
there is a uniform gain. We use this gain twice. First, it yields sparse
matrices whose Poincar\'e constants grow more slowly than the exponent
governing their orders. Then a graph product turns this disparity into
a uniform bound at fixed degree.

The product uses the two equally weighted sets of transitions in the
balanced replacement product of Reingold, Vadhan and Wigderson
\cite[Definition~6.6]{RVW}. For the corresponding nonlinear spectral
comparison, see Mendel and Naor \cite[Theorem~1.13(V) and \S8]{MendelNaor}.
Their Ces\'aro estimates for the two-map nonlinear spectral gap in Hadamard
spaces appear in \cite[Corollary~1.9]{MendelNaorBarycentric}
and \cite[Theorem~3.4]{MendelNaorHadamard}. We give the one-map
lazy-walk estimate with the constants needed here.

Let $\mathcal X_F$ be the class of complete $\operatorname{CAT}(0)$
spectral spaces over $F$. By Theorem~\ref{thm:geometry}, every spectral
wire pseudometric embeds in a member of $\mathcal X_F$ after identifying
points at distance zero. Completeness permits us to take barycenters.
Write $\ell=k+1$ and let $Q_\ell=\{0,1\}^\ell$ be the excluded set of
cube corners with the Euclidean metric. Squared distances on $Q_\ell$
are Hamming distances. A matrix
$G$ has \emph{bound} $\gamma$ if $\mathcal V(x)\le\gamma\mathcal E_G(x)$
for every assignment into every member of $\mathcal X_F$. In a regular
graph we count both directions of an edge and count a loop as one diagonal
arc.

\subsubsection{A strict inequality on the Grassmannian}
Fix a power of two $m>\ell$, put
$W=\mathbb F_2^{2m}$, and let $\mathcal G$ consist of its $m$-dimensional
subspaces. The binary field here indexes the subspaces. Their images
lie in a spectral space over the coefficient field $F$, of arbitrary
characteristic. Given a map of $\mathcal G$ into $X$, let $B_r$ be the mean
squared distance over ordered pairs $(U,V)$ with $\dim(U+V)=m+r$.
Thus $B_1$ averages over adjacent subspaces, which meet in a hyperplane
of each, while $B_m$ averages over complementary subspaces.

\begin{lemma}\label{lem:wire-grassmann-gap}\label{lem:grassmann-gap}
There is an $\eta>0$, independent of $X\in\mathcal X_F$ and the map, such
that $B_m\le(m-\eta)B_1$.
\end{lemma}

\begin{proof}
By \eqref{eq:cur-quad}, the sum of the two squared diagonals of a
quadrilateral is at most the sum of the four squared sides. At
equality its opposite sides have equal lengths, as shown after
\eqref{eq:cur-quad}.

For $r,s>0$ with $r+s\le m$, take spaces $K,R_0,R_1,S_0,S_1$ in direct
sum of respective dimensions $m-r-s,r,r,s,s$. The four spaces
$K+R_i+S_j$ form a quadrilateral indexed by $i,j\in\{0,1\}$. Under a
uniform random automorphism of $W$, each ordered pair is uniform on all
pairs with its given sum dimension. Averaging the quadrilateral inequality
yields
\[
 B_{r+s}\le B_r+B_s,\qquad B_r\le rB_1.
\]
Suppose for a moment that $B_1=1$ and $B_m=m$. Then $B_r=r$ for every
$r$, by $B_m\le B_r+B_{m-r}$ for $0<r<m$. Every quadrilateral just used
is an equality configuration, since the deficits of all its translates are nonnegative
and have average zero.

Fix $0<r<m$ and independent columns $u_1,\ldots,u_r,w_1,\ldots,w_r$ in
$W$, with respective spans $U,U'$. Complete them by $m-r$ further
independent columns of span $S$. The image distance between $S+U$ and
$S+U'$ is independent of this choice of completion. Indeed, replacing one
completion column by a vector outside the current total span preserves the
distance by equality of opposite sides in the quadrilateral of type
$(r,1)$. These replacements connect all completions. To see this, for any
independent tuple of length less than $2m$ with a fixed prescribed
subtuple, one can simulate any independence-preserving change of a single
free column by these replacements. If the desired new column belongs to
the present span, first replace the old one by a vector outside the span.
This puts the desired column outside the new span. Arbitrary
independence-preserving changes connect the tuples by extending them to full
bases and using elementary column operations fixing the prescribed subtuple.
These are generated by additions on the remaining columns over
$\mathbb F_2$ (first modulo the prescribed subspace, then adding its
columns). Record only the changes to the original tuple.

The resulting common distance is also invariant when one of the $2r$
original columns $z$ is changed to a column $z'$ outside their span.
For this move choose the completion to contain $z+z'$ and extend it independently. The
total spans before and after this choice are the same, and neither
$S+U$ nor $S+U'$ changes. Since $2r<2m$, connectivity of the column
moves applies again. Every pair with sum dimension $m+r$ has such a
presentation, so every such squared distance equals $B_r=r$.

Choose $\ell$ disjoint pairs of columns and another $m-\ell$ columns,
all independent. The spaces spanned by the common columns and one choice
from each pair have squared image distances equal to the number of
differing choices, since $\ell<m$. Their images are an isometric copy
of $Q_\ell$.

Finally, if $B_1=0$ then $B_m=0$. If the asserted strict uniform gap
failed, we could rescale a sequence of configurations to have $B_1=1$
and $B_m\to m$. All entries of the finite distance arrays are bounded since
$B_r\le rB_1$ for every $r$. Take a convergent subsequence. Every array
is realizable up to scale over $F$, so the ultralimit procedure of Section~\ref{sec:limits}
realizes its limit up to scale as well. In a corresponding rescaled
geodesic completion this limiting map has $B_1=1$ and $B_m=m$. The
equality argument produces the excluded $Q_\ell$, a contradiction.
\end{proof}

\subsubsection{Sparse matrices on rapidly growing sets}
Decrease $\eta$ if necessary and choose $\varepsilon>0$ so that
$1<\alpha=(m-\eta)(1+\varepsilon)<m$. For every power of two $D\ge2$
and every integer $j\ge0$, we construct a symmetric stochastic matrix on
$D^{m^j}$ vertices with at most $C_jD$ entries in the support of each row
and with bound $C'\alpha^j$. Here $C_j$ may depend on $j$, whereas $C'$
is independent of both $j$ and $D$.

Start with the uniform matrix at size $D$. Given a matrix at size $N$,
let $E=\mathbb F_2^{\log_2 N}$ and take $\Omega=E^m$ as the new vertex
set. Fix an ordered basis in each $U\in\mathcal G$. A uniformly random
linear map $W\to E$ defines a uniform vertex of $\Omega$ at each $U$ by
reading its basis images. For a complementary pair of subspaces the two
vertices are independent. For an adjacent ordered pair $(U,V)$ of
subspaces they are independent conditional on the map on $U\cap V$, and
each conditional vertex is uniform in a fiber of size $N$. For each
endpoint these fibers partition $\Omega$, with uniform fiber weights.
Apply Lemma~\ref{lem:wire-grassmann-gap} to the points assigned to these
vertices and average over the random map. It bounds $\mathcal V$ by
$m-\eta$ times the average of the adjacent-pair expectations.

Fix one ordered adjacency $(U,V)$. For each value of the intersection
map, choose a bijection from the fiber at $U$ to the corresponding fiber
at $V$. Together these bijections give a permutation $\pi$ of $\Omega$.
On each fiber at $U$ place a copy of the preceding matrix, of bound
$\gamma$, and call their union $G_{UV}$. Let $P_\pi$ be the symmetric
average of the permutation matrix and its inverse. For the two conditional
vertices $a,b$, the vertex $\pi^{-1}(b)$ is uniform in the same fiber
as $a$ and independent of $a$ there. Comparing through $\pi^{-1}(b)$
and averaging over the fibers gives
\[
 \mathbb E\,d(x_a,x_b)^2\le
 (1+\varepsilon)\gamma\,\mathcal E_{G_{UV}}(x)
 +(1+1/\varepsilon)\mathcal E_{P_\pi}(x).
\]
The first term uses independent uniform vertices in a single fiber.
For the second, $b$ is uniform on $\Omega$, so the mean squared
distance from $x_{\pi^{-1}(b)}$ to $x_b$ is the permutation-edge energy.

Average the two weighted matrices over all ordered adjacencies and
normalize their sum to be stochastic. The Grassmannian estimate bounds
$\mathcal V$ by the energy of this matrix times
\[
 (m-\eta)\bigl((1+\varepsilon)\gamma+1+1/\varepsilon\bigr)
 =\alpha\gamma+(m-\eta)(1+1/\varepsilon).
\]
Since $\alpha>1$, iteration gives the bound $C'\alpha^j$. There are
only a fixed number of ordered adjacencies, depending on $m$, so each
step increases the support per row by at most a fixed multiple of the
old support plus a constant. At step $j$ the support is at most $C_jD$.

We realize each such matrix $G$ by a regular multigraph at a loss of
at most two in the bound. Choose a power of two $h$ with
$2C_jD\le h\le4C_jD$, increasing $C_j$ to at least one. Form symmetric
arc counts $\lceil h G_{uv}/2\rceil$ and pad each row to total $h$ with
loops. Before padding the row sum is at most $h/2+C_jD\le h$, and the new
energy is at least half the previous energy.

\subsubsection{Bounded-degree amplification}
The extension step replaces each vertex by a graph whose vertices
serve as ports. Within each copy the mean squared edge length controls
dispersion around the barycenter. Between copies a matching of ports realizes several steps of
the old walk. We first estimate the energy of these longer walks.

Let $G$ have bound $\gamma\ge1$,
put $A=(I+G)/2$, and, for even $t\ge2$, put $S=t^{-1}\sum_{h=1}^t A^h$.
Then $S$ has bound at most $C_0(1+\gamma/t)$ for a universal $C_0$.
To prove this, recall that a finite probability distribution on a complete
$\mathrm{CAT}(0)$ space has a barycenter minimizing the mean squared
distance. The midpoint inequality makes each minimizing sequence Cauchy, and
the geodesic squared-distance inequality with the interpolation fraction
tending to zero shows that the mean at $z$ exceeds the minimum by at
least the squared distance from $z$ to the barycenter.

Starting at $x_i^0=x_i$, let $x_i^{g+1}$ be the barycenter of the points
$x_j^g$ under row $i$ of $A$. Set
\[
 E_g=\frac1{|V|}\sum_{i,j}A_{ij}d(x_i^{g+1},x_j^g)^2,
 \qquad D_g=\frac1{|V|}\sum_i d(x_i^g,x_i^0)^2.
\]
Start a length-$h$ walk from the uniform law. Condition on the walk through its
penultimate vertex and apply the barycenter inequality with the initial
point as the test point. Repeat backwards through the walk. Stationarity
gives $\mathcal E_{A^h}(x)\ge D_h+\sum_{g<h}E_g$. Summing in $h$, using
$D_0=0$, yields a $0\le g<t/2$ with
\[
 D_g+(t/2)E_g\le2\mathcal E_S(x).
\]
Comparison through $x_i^{g+1}$ bounds $\mathcal E_G(x^g)\le4E_g$, since
half of row $i$ of $A$ comes from staying at $i$. Comparing $x_i^0,x_j^0$
through $x_i^g,x_j^g$ now gives
\[
 \mathcal V(x)\le 3(2D_g+\mathcal V(x^g))
 \le3(2D_g+4\gamma E_g)
 \le\max\{12,48\gamma/t\}\mathcal E_S(x).
\]

Fix $j$ large and then $D$ large, both specified below. As a graph on a
set of ports take the rounded graph $H$ just constructed on $Q=D^{m^j}$
vertices, with degree $h$ and bound $\beta=\max\{1,2C'\alpha^j\}$.
Put $b=2h$ and $t=m^j/4$, an even power of two for large $j$. We can
ensure $t(2b)^t\mid Q$ because all quantities are powers of two, $2b\le16C_jD$,
and for sufficiently large $D$,
$t(16C_jD)^t\le D^{4t}=Q$.

Given a symmetric $b$-regular graph with normalized matrix $G$, replace
each vertex by a cloud of $Q$ ports. On the new vertex set use $H$ inside
each cloud with probability $1/2$ and otherwise use a port
involution realizing the transitions $S$ between clouds for uniform
ports. Each $A^i$ has denominator dividing $(2b)^t$, so $QS$ is integral.
At cloud $u$ partition the ports into sets of sizes $QS_{uv}$. Pair the parts indexed by $u,v$ and by $v,u$,
using fixed points for self-transitions. Taking the $h$ arcs of $H$ and
$h$ copies of the matching arc at each port gives degree $b$ again.

Let $y_u$ be the barycenter of cloud $u$ and let $V_0$ be the mean squared
distance to one's own barycenter. Write $E_{\rm in},E_{\rm match}$ for the
mean internal and matching energies. By minimization and the bound on $H$,
$V_0\le\beta E_{\rm in}$. The three-segment estimates via a matched pair
and via the two barycenters, respectively, give
\[
 \mathcal E_S(y)\le3(2V_0+E_{\rm match}),\qquad
 \mathcal V(x)\le3(2V_0+\mathcal V(y)).
\]
Since the new energy is $(E_{\rm in}+E_{\rm match})/2$, a new bound is
$C_1\beta(1+C_0(1+\gamma/t))$ with universal $C_1$, whenever the old bound
is $\gamma\ge1$. The coefficient $C_1 C_0\beta/t$ is
$O((\alpha/m)^j)$, so choose $j$ to make it less than $1/2$ and then
choose $D$ as above. Starting with a single vertex with $b$ loops, the
iteration gives graphs on the powers of $Q$, with a uniformly bounded
Poincar\'e constant and degree $b$.

The resulting family has unbounded order, fixed degree and a common
bound for every target in $\mathcal X_F$. This proves
Proposition~\ref{lem:refresh-graph}.
